\section{Related Work}

\subsection{Few-Shot Novel View Synthesis}
Building upon the foundational NeRF framework, recent methodologies have emerged to address the challenges of few-shot novel view synthesis. Early approaches in this domain typically rely on strong regularization priors. For instance, RegNeRF~\cite{niemeyer2022regnerf} regularizes geometry by enforcing smoothness constraints on unobserved viewpoints, while FreeNeRF~\cite{yang2023freenerf} employs frequency regularization to prevent HF artifacts during the early stages of training. Other methods, such as SparseNeRF~\cite{roessle2022sparsenerf} and DietNeRF~\cite{jain2021putting}, incorporate auxiliary supervision to guide reconstruction. Specifically, SparseNeRF utilizes depth priors, while DietNeRF leverages semantic consistency losses to guide geometry in occluded regions.

Recent advancements have focused on robust geometric adaptation without heavy external priors. FrugalNeRF~\cite{lin2025frugalnerf} introduces a cross-scale sharing scheme to maximize information utility from limited pixels. Furthermore, addressing the practical reality of agricultural and robotic data, methods like SPARF~\cite{truong2023sparf} have extended few-shot capabilities to handle noisy camera poses, jointly refining extrinsic parameters and scene geometry to prevent drift in uncontrolled environments.

\subsection{3DGS and Sparse-View Extensions}

3DGS substantially reduces the training latency of NeRF-based methods through efficient differentiable rasterization. Despite this advantage, 3DGS exhibits a characteristic failure mode in sparse-view or few-shot regimes. In such settings, the optimization process tends to over-reconstruct HF details, particularly edges and fine textures, in the observed training views. Due to limited multiview constraints, this behavior leads to poor generalization across viewpoints and results in degraded global geometric coherence and weakened structural consistency in unobserved regions~\cite{nguyen2025dwtnerf}.

To address these limitations, several extensions have been proposed to improve the robustness of 3DGS under sparse supervision. Methods such as FSGS~\cite{zhu2024fsgs} and PGDGS~\cite{huang2025pgdgs} employ adaptive and progressive densification strategies that incrementally populate underconstrained regions of the scene, thereby reducing geometric sparsity and improving reconstruction stability. Moreover, DNGaussian~\cite{li2024dngaussian} and SCGaussian~\cite{peng2024structure} focus on regularizing scene geometry through explicit depth constraints and structural consistency priors. By suppressing spurious or unstable Gaussian primitives, these approaches enhance geometric plausibility and reduce reconstruction artifacts in sparsely observed areas. While effective, such methods primarily rely on geometric supervision and do not explicitly regulate the spectral distribution of reconstructed content, leaving frequency-domain inconsistencies insufficiently constrained.


\subsection{Frequency-Aware Supervision}

Frequency-domain analysis has been adopted in neural rendering as an effective means to separate global structural information from fine-scale texture. LF components primarily encode smooth geometry and large-scale scene structure, whereas HF components correspond to edges and detailed appearance variations. By explicitly regulating these components, frequency-aware supervision provides a principled mechanism for stabilizing optimization under limited viewpoint coverage.WaveNeRF~\cite{xu2023wavenerf} and DWT-NeRF~\cite{nguyen2025dwtnerf} use wavelet guidance to improve NeRF training. This reduces errors and sharpens edges, especially when there are few camera views. Similarly, DWT-GS~\cite{nguyen2025dwtgsrethinkingfrequencyregularization} applies this concept to Gaussian Splatting to remove HF noise. However, these methods generally apply rules to the whole image at once. They fail to distinguish between preserving the main structure and refining fine details."

Building on these observations, our approach introduces a dual-branch frequency-aware supervision strategy that is directly integrated into the 3DGS rasterization pipeline. A global frequency branch enforces LF consistency to preserve overall geometric structure, while a complementary patch-wise branch selectively refines HF components to recover local details without inducing overfitting. 


\subsection{Multispectral and Agricultural 3D Reconstruction}

Multispectral and hyperspectral imaging capture reflectance information beyond the visible spectrum, enabling critical applications in agriculture~\cite{karukayil20253d}, remote sensing~\cite{ziemann2025new,nia2025neighborhood}, and material analysis~\cite{klein2023physics}. By providing wavelength-dependent measurements, these modalities support robust characterization of vegetation health, structural properties, and material composition that cannot be inferred from RGB imagery alone. Several NeRF-based extensions have been proposed to model spectral radiance fields across multiple wavelengths. Methods such as HS-NeRF~\cite{chen2024hyperspectral}, SpectralNeRF~\cite{li2024spectralnerf}, and Spec-NeRF~\cite{li2024spec} explicitly reconstruct spectral reflectance by conditioning radiance fields on wavelength information. More recently, HyperGS~\cite{thirgood2025hypergs} adapts Gaussian Splatting to hyperspectral scenes, demonstrating the feasibility of point-based spectral rendering. However, these approaches generally assume dense multi-view acquisition and rely on high-cost sensing hardware, limiting their applicability in real-world agricultural settings.

In practice, agricultural imaging is often performed under few-view, spectrally heterogeneous conditions, with additional challenges arising from varying illumination, plant self-occlusions, and complex scene geometry. To better reflect these constraints, we introduce a controlled multispectral greenhouse dataset capturing Red, Green, Red Edge, and NIR channels, along with a unified few-shot benchmarking package. This dataset enables systematic evaluation of frequency-aware multispectral 3D reconstruction methods under realistic agricultural imaging conditions.

